\section{Regrouping the exclusion correction}
\label{sec:forced-exclusion}
The regular gas in \eqref{eq:forest-regular-separated} has been
extracted, but the forced family still carries its exclusion factor.
We now put this factor into local forced activities. The factorials
and signs of the repeated logarithmic terms are retained throughout.

Fix a coarse configuration $V$. Continue to write $k_j$ for a forced
activity with support $Y_j$ and nonempty inventory $I_j\subseteq Y_j$,
and $c_\Gamma$ for a regular logarithmic cluster term with nonempty
support $X_\Gamma$. We have summable majorants
\[
 |k_j|\le p_j,\qquad |c_\Gamma|\le q_\Gamma,\qquad
                  \sum_jp_j+\sum_\Gamma q_\Gamma<\infty.
\]
These follow on a finite lattice from the previously purchased rooted
bounds. Here $q_\Gamma$ is specifically the absolute Ursell coefficient
divided by the tuple factorial, multiplied by the purchased regular
constituent prices. Thus the rooted cluster estimates apply to this
majorant itself, not merely to $|c_\Gamma|$.
All component addresses, support sets and eligibility tests
are fixed before their values at $V$ are evaluated.

For a finite admissible forced family $A$, the exclusion correction
has the absolutely convergent expansion
\begin{equation}\label{eq:exclusion-multiplicities}
 e^{-P_{Y_A}}=
 \sum_{\substack{n:\ \text{finite support in the cluster alphabet}\\
                n_\Gamma>0\Rightarrow X_\Gamma\cap Y_A\ne\varnothing}}
               \prod_{\Gamma:n_\Gamma>0}
                                   \frac{(-c_\Gamma)^{n_\Gamma}}{n_\Gamma!},
 \qquad n_\Gamma\in\mathbb N.
\end{equation}
For a finite cluster alphabet this is the product of the ordinary
exponential series. For the whole alphabet, the sum of the absolute
products is at most $\exp(\sum_\Gamma q_\Gamma)$; increasing finite
subsets and dominated convergence prove the identity. Moreover the
sum after multiplication by $\prod_{j\in A}k_j$ and summation over
$A$ is absolutely dominated by
\begin{equation}\label{eq:exclusion-absolute}
                 \prod_j(1+p_j)\exp\left(\sum_\Gamma q_\Gamma\right).
\end{equation}
The restriction that every retained cluster meet $Y_A$ is essential:
the correction deletes exactly those regular clusters.

A finite decorated family consists of a finite set of forced labels
and a finite set of cluster labels, with a positive integer multiplicity
attached to each cluster label. Give a forced label its support $Y_j$
and multiplicity one, and a cluster label its support $X_\Gamma$.
Its graph has one vertex per distinct label, with edges for intersecting
supports. A new forced label $D$ is a connected decorated family
such that the old forced supports are pairwise disjoint and every
cluster support meets the union of the old forced supports in $D$.
Write $Z_D$ for its union support and
$I_D=\bigcup_{j\in D\ {\rm forced}}I_j$ for its inventory. In
particular $D$ contains a forced label and $I_D\ne\varnothing$.

Define its value by
\begin{equation}\label{eq:exclusion-new-value}
 \widetilde k_D(V)=
  \mathbf1_{\{\operatorname{Flags}(V)\cap Z_D=I_D\}}
  \mathbf1_{\{\text{the old local eligibility tests hold}\}}
       \prod_{j\in D\ {\rm forced}}k_j(V)
       \prod_{\Gamma\in D\ {\rm cluster}}
                           \frac{(-c_\Gamma(V))^{n_\Gamma}}{n_\Gamma!}.
\end{equation}
The cluster label itself already includes its original tuple
factorial in $c_\Gamma$. The additional $n_\Gamma!$ belongs to the
exponential in \eqref{eq:exclusion-multiplicities}; these are two
different factorials.

Consider a nonzero term in the expansion for $(A,n)$. Split its
support graph into connected components. Every retained cluster
meets a forced support in $A$, and that forced label is in its
component. Thus each component is a new forced label as just
defined. Its inventory equals the flags inside its union support:
a flag there belongs to some old inventory because $A$ covers all
flags, and the corresponding old forced support meets this component,
so belongs to it. The reverse inclusion follows from the original
inventory condition. Old eligibility tests hold for each nonzero
factor. Hence the product of \eqref{eq:exclusion-new-value} over the
components equals the original term, with every sign and factorial.

Conversely, flatten a finite compatible family of new labels whose
inventories cover all flags and whose product is nonzero. Its old
forced labels are pairwise compatible, both within and between the
new labels. They satisfy all eligibility and exact flag-coverage
conditions. A cluster label cannot occur in two new components:
its nonempty support would intersect both of their supposedly
disjoint union supports. Its multiplicity is therefore recovered
unambiguously. Each retained cluster meets a forced support in its
own component, as required in \eqref{eq:exclusion-multiplicities}.
The connected components of the flattened graph are exactly the
original new labels. Flattening and taking components are inverse.

This weight-preserving bijection and \eqref{eq:exclusion-absolute}
give
\begin{equation}\label{eq:exclusion-output-gas}
 (\mathcal BR)(V)=\mathscr O(V)
   \exp\left(\sum_\Gamma c_\Gamma(V)\right)
   \sum_{\substack{\mathcal D\ {\rm finite,\ compatible}\\
                    \bigcup_{D\in\mathcal D}I_D=\operatorname{Flags}(V)}}
                                      \prod_{D\in\mathcal D}\widetilde k_D(V).
\end{equation}
If there are no flags, no nonzero forced component is possible and
the empty family contributes one. If there are flags, every nonzero
component has a nonempty inventory. This is the required separation
into a regular exponential and a forced gas.

The new values are measurable and gauge invariant: each factor and
eligibility predicate is, and the products are finite. They use only
the final coarse field and its flags on $Z_D$, by locality of the
constituents and of the inventory test. Under a signed lattice
transformation the old labels and the cluster occurrence tuples are
permuted, with their multiplicities unchanged. Thus the value,
support and inventory transform together. A common countable
enumeration may then be chosen for the new labels, assigning zero
to any unused addresses. Such an enumeration only reindexes the
absolutely convergent gas and changes none of its values or support
conditions.

For later bounds the same construction has the explicit nonnegative
majorant
\begin{equation}\label{eq:exclusion-output-price}
           \widetilde p_D=
              \prod_{j\in D\ {\rm forced}}p_j
              \prod_{\Gamma\in D\ {\rm cluster}}
                                      \frac{q_\Gamma^{n_\Gamma}}{n_\Gamma!}.
\end{equation}
Taking absolute values in \eqref{eq:exclusion-new-value} proves the
bound by \eqref{eq:exclusion-output-price}. We next prove the uniform
rooted estimate for these majorants.

\subsection{Keeping one compulsory gain through the second expansion}
Let $a\ge0$ be the desired support exponent and $\varepsilon>0$ a
reserve. Suppose the regular constituent prices obey the cluster
criterion at support exponent $a+\varepsilon$, rooted allowance
$u>0$ and tail ratio $1/2$. Suppose also that the old forced prices
obey
\[
 \sup_x\sum_{j:x\in Y_j}p_j e^{(a+\varepsilon)c_B(Y_j)}
                                                   \le u e^{-M},
                      \qquad M\ge0.
\]
The regular cluster bound then gives
\[
 \sup_x\sum_{\Gamma:x\in X_\Gamma}
                   q_\Gamma e^{(a+\varepsilon)c_B(X_\Gamma)}\le2u.
\]
Set $\eta=20u$ and assume
\begin{equation}\label{eq:exclusion-second-smallness}
       \eta\le\tfrac12,\qquad
       \frac{4b_B\eta}{\varepsilon}\le\tfrac12,
                              \qquad b_B=(2B+1)^4.
\end{equation}
The second inequality is the explicit connected-tree degree bound
of Section~\ref{sec:cluster-expansion}.

We prove that, for every site $x$ and integer $m\ge0$,
\begin{equation}\label{eq:exclusion-rooted-final}
 \sum_{\substack{D:\,x\in Z_D\\ |D|_{\rm copies}\ge m}}
              \widetilde p_D e^{a c_B(Z_D)}
                         \le 2^{-m}e^{-M}\,2\eta.
\end{equation}
For an upper bound forget the old compatibility tests, allow forced
labels to repeat as well, and retain only connectedness and the
presence of at least one forced label. On the resulting combined
alphabet put $w_j=p_j$, $w_\Gamma=q_\Gamma$, and give each forced
occurrence a mark $M$, each regular occurrence a mark zero.
Inflate a price to
\[
             v_i=2w_i e^{e_i},\qquad e_j=M,\quad e_\Gamma=0.
\]
The rooted bound for $2v_i$ at exponent $a+\varepsilon$ is at most
$4e^Mue^{-M}+4(2u)=12u\le\eta$. Thus these prices meet the
distinct connected-set criterion with allowance $\eta$ and tail
ratio $1/2$, by \eqref{eq:exclusion-second-smallness}.

To include multiplicities, first charge the desired support weight
to each occurrence by setting $v_i'=v_i e^{a c_B(X_i)}$, with $X_i$
denoting either kind of support. Since $X_i\ne\varnothing$, choose
a site in $X_i$ and use the single term of the rooted bound for
$2v_i$ to obtain $v_i'\le\eta\le1/2$. For a fixed distinct set of
labels $S$, summing all positive multiplicities gives
\[
 \sum_{(n_i\ge1)_{i\in S}}\prod_{i\in S}\frac{(v_i')^{n_i}}{n_i!}
       \le\prod_{i\in S}\sum_{n\ge1}v_i'(1/2)^{n-1}
       =\prod_{i\in S}2v_i'.
\]
Here we used $n!\ge1$ and $v_i'\le1/2$. Connectedness depends only
on the distinct supports, so the distinct connected-set bound now
shows that the sum over all rooted connected multiplicity families
is at most $\eta/(1-1/2)=2\eta$. Cover subadditivity on a connected
set gives
$c_B(\bigcup_{i\in S}X_i)\le\sum_{i\in S}c_B(X_i)
\le\sum_i n_i c_B(X_i)$, justifying the charge per occurrence.

Finally the inflated product differs from the original one by
$2^{\sum_i n_i}\exp(\sum_i n_i e_i)$. For a retained family these
two factors are at least $2^m$ and $e^M$, respectively. Extracting
them proves \eqref{eq:exclusion-rooted-final}. Restricting back to
the actual new labels only decreases the nonnegative sum. This
argument preserves a compulsory gain even after arbitrary repeated
regular corrections have been attached.

In the selected scale estimates, $u(x)$ has a fixed negative upper
exponent and $B(x)$ has upper exponent zero. Both inequalities in
\eqref{eq:exclusion-second-smallness}, and the first regular degree
bound with its own fixed positive reserve, therefore hold eventually.
The compulsory mark is
\[
                M(x)=\frac{\mu}{4}\operatorname{fineMonitor}(x)^2,
                            \qquad \mu>0.
\]
Thus \eqref{eq:exclusion-rooted-final} is precisely
$40u(x)2^{-m}e^{-M(x)}$. It pays every fixed outgoing gain
$\exp(-A_0 P(x)^\delta)$ with $A_0\ge0$ and $0\le\delta<2$,
where $P(x)=\exp((\log\log x)^2)$ is the unshifted monitor.
Indeed $P(x)\le\operatorname{fineMonitor}(x)$ eventually, the ratio
$\operatorname{fineMonitor}(x)^\delta/
\operatorname{fineMonitor}(x)^2$ tends to zero, and $40u(x)2^{-m}\le1$
eventually. This proves the outgoing forced allowance from the
incoming regular/forced prices, with the same scale and history.
